\section{Introduction}\label{sec:intro}
\pagenumbering{arabic} \setcounter{page}{1}

In Singapore, a compact grid serves approximately 5.9 million residents under a tropical climate that sustains year-round cooling demand; national electricity consumption exceeded 55\,TWh in 2024 \parencite{ema2025statistics}. Because air-conditioning is the largest end-use of electricity in buildings, demand is highly sensitive to small variations in temperature and humidity \parencite{zhao2012review}, with holidays, weekends, and monsoon seasons adding further variability. Accurate short-term forecasting is therefore essential for reliable grid operation and cost-effective generation scheduling.

\textbf Classical methods (ARIMA, exponential smoothing, regression) assume linear, stationary relationships and struggle with the non-linear interactions among weather, calendar, and autoregressive demand dynamics \parencite{almusaylh2018forecasting}. Tree-based machine learning models capture these interactions well but are sensitive to hyperparameters, and manual tuning is expensive and often sub-optimal. Zhang et al.\ \parencite{zhang2023catboostppso} showed that CatBoost paired with Phasor Particle Swarm Optimisation (PPSO) outperforms five other meta-heuristics---but on hourly data from a \emph{single residential customer}, and with tuning applied only to the favoured model. Whether the approach generalises to \emph{daily national-level} demand in a tropical setting, and whether its advantage survives when \emph{all} competing models receive the same tuning budget, remain open questions.

\textbf This project (i) builds a daily Singapore dataset merging EMA demand with Open-Meteo weather data (February 2023 -- December 2025); (ii) engineers 16 features across weather, derived-comfort, calendar, and autoregressive categories; (iii) compares eight model configurations---a Lag-1 baseline, Linear Regression, Random Forest, XGBoost, and CatBoost, plus PPSO-tuned variants of each tree model under an \emph{identical} budget---on six metrics; and (iv) deploys the best model through a web portal (FastAPI + React) for non-technical users. Intra-day resolution, sector decomposition, and deep learning models are out of scope.
